
% --- CAPÍTULO 5 ---
\chapter{Conclusão}
Este trabalho de conclusão de curso abordou o desafio da predição de diabetes tipo 2, um problema de saúde pública de crescente relevância e prevalência global. O objetivo central foi desenvolver e avaliar modelos de redes neurais artificiais capazes de funcionar como uma ferramenta de apoio ao diagnóstico precoce, utilizando dados clínicos e comportamentais acessíveis.

O problema foi tratado através da construção de um pipeline de \textit{Machine Learning} robusto, aplicado a um \textit{dataset} público de 100.000 registros. Um desafio metodológico central foi o significativo desbalanceamento de classes do \textit{dataset} (92% não diabéticos e 8% diabéticos). Para mitigar o risco de um modelo enviesado para a classe majoritária, foi implementada a técnica de \textit{data augmentation} SMOTE (Synthetic Minority Over-sampling Technique), permitindo um treinamento mais equilibrado e focado na correta identificação dos casos positivos.

Foram implementadas e comparadas duas arquiteturas de redes neurais, uma \textit{Multi-Layer Perceptron} (MLP) e uma Rede Neural Convolucional (CNN) 1D , além de modelos clássicos de \textit{Machine Learning} (Random Forest e Gradient Boosting) para estabelecer um \textit{baseline}. A otimização bayesiana foi utilizada para o ajuste fino dos hiperparâmetros das redes neurais, garantindo uma exploração eficiente do espaço de busca. Os resultados obtidos foram promissores: os modelos de redes neurais alcançaram acurácia elevada, com o MLP atingindo 89,11% e a CNN 88,96%  após o retreinamento. Mais importante que a acurácia, o foco na otimização de métricas como \textit{recall} (onde o MLP atingiu 91,98% ) e a análise da curva PR-AUC  demonstrou a capacidade dos modelos em minimizar os Falsos Negativos. Esta foi uma decisão de projeto deliberada, dado que falhar em diagnosticar um paciente doente (Falso Negativo) representa o cenário de maior risco clínico.

A contribuição deste trabalho pessoalmente, a execução deste projeto exigiu uma imersão profunda nas complexidades da aplicação de inteligência artificial em um contexto de alto impacto social. O desafio de tratar dados desbalanceados , implementar e otimizar arquiteturas de \textit{Deep Learning} , e interpretar os resultados sob a ótica da responsabilidade ética (priorizando o \textit{recall} para minimizar o risco clínico ), solidificou competências técnicas em ciência de dados, engenharia de \textit{Machine Learning} e pensamento crítico.

Finalmente, este trabalho cumpre seus objetivos ao entregar modelos funcionais e ao validar a hipótese de que redes neurais são ferramentas eficazes para esta tarefa. Conforme detalhado na seção de Trabalhos Futuros , o pipeline desenvolvido serve como alicerce para a criação de uma plataforma interativa e a validação com \textit{datasets} mais abrangentes, reforçando o potencial de aplicação prática da solução.